\section*{Phase 1: Data Science Project Proposals}

\subsection*{Project 1: Housing Amenities Impact on Apartment Listing Prices}

\subsubsection*{Business Issue and Overview}

Real estate investors and property managers lack quantitative evidence on which amenities drive rental pricing. Understanding amenity importance enables strategic property improvements and competitive pricing. This project analyzes which amenities (parking, laundry, pet policies, fitness centers, etc.) most significantly influence rental prices across U.S. markets.

\subsubsection*{Business Questions}

\begin{itemize} [itemsep = -9pt]
    \item Which amenities have strongest impact on rental pricing?
    \item How does amenity importance vary across geographic regions?
    \item Which amenity upgrades yield the highest price premium relative to cost?
    \item Do certain amenity bundles (e.g., parking \& laundry) have compounding effects on price?
\end{itemize}

\subsubsection*{Data Needs}

Major entities: Rental listings (price, location, bedrooms, bathrooms, square footage), amenities (binary/categorical flags), geographic clusters (region, state), listing metadata (date listed, availability).

\subsubsection*{Reason}

This project addresses a real business problem faced by property investors and managers daily. Understanding which amenities justify investment is directly actionable, it informs renovation priorities and pricing strategy. Personally, this project deepens familiarity with feature importance techniques and regression modeling, skills broadly applicable across industries.

\newpage

\subsection*{Project 2: Student Course Success Prediction}

\subsubsection*{Business Issue and Overview}

Educational institutions face resource constraints and must identify at-risk students early. This project builds a predictive model to identify students likely to struggle or fail in courses, enabling early intervention programs. Predictions target first-semester students to maximize intervention effectiveness.

\subsubsection*{Business Questions}

\begin{itemize} [itemsep = -9pt]
    \item Which student characteristics predict course success?
    \item When in the semester should interventions occur?
    \item What baseline GPA threshold indicates risk?
    \item How accurately can we identify at-risk students?
\end{itemize}

\subsubsection*{Data Needs}

Major entities: Student demographics (age, prior GPA, major, course history), course characteristics (difficulty, subject, instructor), performance metrics (grades, attendance, engagement).

\subsubsection*{Reason}

Addresses critical institutional challenge: student retention and success. Interest in educational technology and predictive analytics applied to real student outcomes. Classification modeling provides interpretable insights; understanding why students succeed enables actionable intervention strategies. Valuable for education stakeholders and demonstrates practical impact of data science on human outcomes.

%\subsection*{Recommended Project}

%\textbf{Project 1 (Housing Amenities)} is recommended. It aligns with your prior work, has established data availability (Craigslist listings), and delivers concrete business value for real estate stakeholders. The feature importance analysis directly answers investor questions about amenity ROI.
